\section{Related work}
\label{sec:related}

\paragraph{General time-series benchmarks.} Monash \citep{godahewa2021monash}, GIFT-Eval \citep{aksu2024gifteval},
TFB \citep{qiu2024tfb} and TSFM-Bench \citep{li2025tsfmbench} evaluate across many domains and report aggregate
ranks. Weather appears as one domain among dozens, drawn mostly from mid-latitude sources, and none stratifies by
climatic regime or geography. TIME \citep{qiao2026time} adds a strict leakage-free protocol over 50 fresh datasets
and 12 foundation models, and likewise applies no geographic or climatic strata. Recent work has begun to document
that aggregate metrics conceal regime-dependent failure \citep{wang2026regime} and that benchmark construction
itself is a source of leakage \citep{meyer2026leakage}; our design follows both warnings.

\paragraph{Weather benchmarks.} WeatherBench 2 \citep{rasp2024weatherbench2} and ChaosBench
\citep{nathaniel2024chaosbench} score against reanalysis; ExtremeWeatherBench \citep{mcgovern2026ewb} scores
against event catalogues; WeatherReal \citep{jin2024weatherreal} and WEATHER-5K \citep{han2026weather5k} score against
station observations. WEATHER-5K is the closest in spirit (thousands of global stations, hourly), but carries
neither precipitation nor humidity, which are the variables at issue here, and reports global aggregates rather
than climatic strata. None of these evaluates time-series foundation models on a tropical station network.

\paragraph{Geographic disparity in AI forecasting.} SAFE \citep{masi2025safe} stratifies AI weather-prediction
models by country, subregion and income group, finding disparities across every attribute, but works on ERA5 at
\ang{1.5} resolution for upper-air fields only. \citet{mozaffari2026inequality} argue qualitatively that AI risks
amplifying a North--South divide without supplying measurements. \citet{rowell2026brazil} stratify GraphCast by
climatic sub-region over Brazil against operational IFS analysis for \(T_{850}\), \(Q_{850}\) and \(Z_{500}\),
excluding precipitation. Ours is the corresponding measurement for time-series foundation models, on station
observations, including rainfall.

\paragraph{South Asian monsoon and AI models.} \citet{gupta2026mausam} assess seven global AI weather-prediction
models against South Asian observations during the monsoon, reporting errors 15--45\% larger than against
reanalysis, the precedent for our observation-versus-reanalysis contrast, transferred here to a different model
class. \citet{luitel2026ism} evaluate GraphCast over the Indian summer monsoon against ERA5 and IMERG, and
\citet{aitken2026monsoon} build decision-oriented probabilistic onset forecasts for Indian farmers from AI weather
models. All three concern gridded NWP-style emulators and onset or seasonal-mean skill rather than zero-shot
station forecasting.

\paragraph{Precipitation-specific evaluation.} A recent benchmark evaluates 17 supervised architectures for 0--6\,hour
GNSS-based precipitation nowcasting \citep{zhang2025rainfallbench}. The task, horizon and model class all differ from
ours, but it is the closest existing precipitation-focused time-series benchmark.

\paragraph{Bangladesh forecasting.} The applied literature is substantial but uses per-station supervised models
and releases no shared tasks: \citet{joy2026bangladesh} is representative, training attention-enhanced LSTMs for
temperature and rainfall. No published work evaluates pretrained foundation models on Bangladeshi station data, and
no Bangladesh weather benchmark with fixed splits and a leaderboard exists.

\paragraph{Position.} The intersection is, to our knowledge, unoccupied: current time-series foundation models, on
tropical station observations including precipitation, stratified by monsoon regime, with a matched temperate
control and a significance protocol. We claim novelty only at that intersection, and cite the neighbours above for
each component individually.
